% Generated draft: requires mathematical and source review.
\documentclass[11pt]{article}
\usepackage[margin=25mm]{geometry}
\usepackage{fontspec}
\setmainfont{Latin Modern Roman}
\usepackage{amsmath,amssymb,mathtools,unicode-math}
\setmathfont{Latin Modern Math}
\usepackage{enumitem,needspace,xurl,hyperref}
\setlength{\parindent}{0pt}
\setlength{\parskip}{5pt}
\setlength{\emergencystretch}{3em}
\newcommand{\N}{\mathbb{N}}
\newcommand{\Z}{\mathbb{Z}}
\newcommand{\Q}{\mathbb{Q}}
\newcommand{\R}{\mathbb{R}}
\newcommand{\C}{\mathbb{C}}
\newcommand{\F}{\mathbb{F}}
\newcommand{\A}{\mathbb{A}}
\newcommand{\PP}{\mathbb{P}}
\newcommand{\E}{\mathbb{E}}
\newcommand{\eps}{\varepsilon}
\newcommand{\dd}{\,\mathrm{d}}
\begin{document}
\section*{UnsolvedMath ID: 30002578}
\textbf{Source problem code:} catalogue entry
\subsection{Definitions and mathematical statement}
Let \(k\) denote a ground field. The source does not specify \(k\), its characteristic, or dimensionality assumptions on the Hopf algebras; these are unresolved parameters of the questions below. A Hopf algebra \(H\) has multiplication, unit, comultiplication \(\Delta\), counit \(\varepsilon\), and antipode \(S\). A Hopf subalgebra is a subalgebra that is also a subcoalgebra and is stable under the antipode. A coalgebra is pointed if every simple subcoalgebra is one-dimensional. A pointed Hopf subalgebra \(P\subseteq H\) is maximal if there is no pointed Hopf subalgebra strictly between \(P\) and \(H\). Maximality does not itself specify a unique greatest pointed Hopf subalgebra.

The supplied record combines three questions, with the following unresolved qualifications.
\begin{enumerate}
\item \textbf{Diagram folding of Nichols algebras.} In the usual construction, for an object \(V\) in a suitable braided setting, the Nichols algebra \(\mathcal B(V)\) is the quotient of the braided tensor algebra \(T(V)\) by its largest homogeneous braided Hopf ideal having zero intersection with degrees zero and one. The source asks: what representation-theoretic construction or phenomenon corresponds to diagram folding of Nichols algebras under the associated representation-theoretic correspondence?

The braided setting, the objects \(V\), the diagrams, the allowed folding operations, and the representation-theoretic correspondence are not identified. Consequently, this question is retained as a request for an interpretation; a particular correspondence or compatibility identity cannot be specified from the supplied record.

\item \textbf{Splitting over a maximal pointed Hopf subalgebra.} Let \(H\) be a Hopf algebra and let \(P\subseteq H\) be the subalgebra denoted \(H'\) in this part of the source, described there as its maximal pointed Hopf subalgebra. A Hopf ideal \(I\subseteq H\) is a two-sided ideal satisfying
\[
\Delta(I)\subseteq I\otimes_k H+H\otimes_k I,
\qquad \varepsilon(I)=0,
\qquad S(I)\subseteq I.
\]
The source asks whether there exists such an ideal with
\[
H=P\oplus I,
\]
and which algebraic and homological relations connect \(H\) and \(P\).

The meaning of the direct sum is not specified. If it means an internal direct sum of \(k\)-vector spaces, the requested condition is precisely
\[
\exists I\subseteq H:\quad I\text{ is a Hopf ideal},\qquad
H=P+I,\qquad P\cap I=\{0\}.
\]
This is a conditional reading, not a resolution of the source ambiguity. The source also does not say whether \(P\) is a greatest pointed Hopf subalgebra, an arbitrary maximal one, or a distinguished choice, and it does not identify the algebraic or homological invariants to be related. Thus the quantification over \(P\) and the scope of the requested relations remain unresolved.

\item \textbf{Frobenius--Schur indicators and gauge equivalence.} Let \(H_1,H_2\) denote the two Hopf algebras in the final source question; no inclusion between them is assumed. For a specified theory of indicators, write formally
\[
\mathcal V(H)=\{\nu_n(M): n\in\mathcal N_H,\ M\in\mathcal M_H\},
\]
where \(\mathcal M_H\) is the intended class of modules, \(\mathcal N_H\) is the intended set of indicator indices, and \(\nu_n(M)\) is the corresponding Frobenius--Schur indicator whenever defined. These choices, the hypotheses ensuring definition, and a common value domain for comparing the two sets are not supplied.

The requested implication is
\[
\mathcal V(H_1)=\mathcal V(H_2)
\quad\Longrightarrow\quad
\text{the module categories of }H_1\text{ and }H_2
\text{ are gauge equivalent}.
\]
The source asks about equality of sets of values; it supplies no module-by-module matching, preservation of indicator indices, or multiplicities. It does not specify whether the module categories consist of all modules or finite-dimensional modules. It also does not define gauge equivalence or its relationship here to tensor equivalence or twisting of Hopf algebras. These omissions prevent a single fully quantified formulation of the implication.
\end{enumerate}
All three questions are retained as statement-only research questions. The supplied open-status assessment and verification labels are curation metadata, not mathematical conclusions established in this draft.
\subsection{Short English statement}
The record asks for an interpretation of Nichols algebra diagram folding, a possible Hopf-ideal complement to a maximal pointed Hopf subalgebra and associated structural relations, and a converse from equality of Frobenius--Schur indicator-value sets to gauge equivalence of module categories. Essential hypotheses and several intended constructions are unspecified, so the draft records these gaps explicitly. The supplied open label is not independently verified.
\subsection{Sources}
\begin{itemize}
\item Supplied local UnsolvedMath registry record, permanent ID \(30002578\), title: \emph{Diagram Folding and Structural Equivalence of Hopf Algebras}; collection: Oberwolfach Reports Open Problems; exact source problem code: \texttt{catalogue entry}. The record supplies proposed year \(2014\), no named proposer, status \texttt{open}, and statement status \texttt{corrected\_verified}. Its verification note says that mathematical content was unchanged while notation, typography, and wording were normalized into a self-contained publication form. These are supplied registry claims, not independent verification.
\item Supplied report citation: \emph{Mini-Workshop: Infinite Dimensional Hopf Algebras (2015)}, Oberwolfach Reports, DOI \texttt{10.4171/owr/2014/20}, pp.\ 1111--1137. Source URL and literature-source link labelled ``Oberwolfach open-questions report'': \url{https://doi.org/10.4171/owr/2014/20}. The supplied background attributes the questions to mathematical workshop open-problem discussions documented in Oberwolfach Reports. The citation's \(2015\) designation, the DOI's \(2014\) component, and the registry's proposed year \(2014\) are preserved without resolving their relationship.
\item Supplied literature-source link: \emph{Total Frobenius--Schur indicators}, \url{https://arxiv.org/abs/1212.1435}. The supplied record does not identify a theorem or passage in this reference that precisely formulates or resolves the final question.
\item Literature assessment checked at \texttt{2026-08-21}: this exact date and the following assessment are supplied dated curation, not independently verified. The record labels the problem open and reports substantial invariant-theoretic theory for diagram folding and indicator values, while stating that the requested interpretation, splitting and homological relations for maximal pointed subalgebras, and converse from indicator sets to gauge equivalence are not established in general. The supplied background describes this as dated literature triage derived from OWR curation and requests independent verification of time-sensitive status and cited claims. No source was visited for this draft.
\end{itemize}
\end{document}
